% §6 Path-Based Tree Encodings
%
% Presents Construction B: the universal path-based method.
% Shows how it works, why it always satisfies the CAR,
% and how it recovers all known encodings.

The path-based tree encoding~\cite{jiang2020, miller2023bonsai}
provides an entirely different route from trees to Pauli strings ---
one that works for \emph{every} tree, not just stars.

\subsection{Encoding trees}
\label{sec:encoding-trees}

\begin{definition}[Encoding tree]
\label{def:encoding-tree}
An \emph{encoding tree} for $n$ fermionic modes is a rooted tree $T$
with internal nodes, where each node has at most three descending
connections, each labelled by a distinct Pauli operator from $\{X, Y,
Z\}$.  Each connection is either an \emph{edge} (to a child node) or a
\emph{leg} (a dangling terminal link).  The total number of legs is
$2n$, paired into $n$ mode assignments.
\end{definition}

The branching constraint (at most 3 children per node) comes directly
from the Pauli algebra: there are exactly three non-identity
single-qubit Paulis ($X$, $Y$, $Z$), and each link label must be
distinct among siblings to ensure anticommutation.

\subsection{The path-based construction (Construction B)}
\label{sec:construction-b}

For mode~$j$, let $\ell_j^{(c)}$ and $\ell_j^{(d)}$ be the paired
legs assigned to that mode.  The Majorana operator for each leg is the
Pauli string obtained by collecting the link labels along the
root-to-leg path:
\begin{equation}
  S(\ell) = \bigotimes_{v \in \mathrm{path}(\ell)}
    \text{label}(v \to \text{next}) \;\otimes\;
    \bigotimes_{v \notin \mathrm{path}(\ell)} I.
  \label{eq:path-pauli}
\end{equation}

Each qubit in the register corresponds to an internal node of the
tree.  At each node on the root-to-leg path, the Pauli label of the
outgoing edge/leg becomes the operator on that qubit.  All nodes
\emph{not} on the path receive identity.

The Majorana operators for mode~$j$ are:
\begin{equation}
  c_j = S(\ell_j^{(c)}), \qquad d_j = S(\ell_j^{(d)}).
\end{equation}
The ladder operators follow from~\eqref{eq:encoded-ladder}.

\subsection{Correctness: why Construction B always works}
\label{sec:tree-car-proof}

\begin{theorem}[Tree encoding correctness]
\label{thm:tree-car}
For any encoding tree $T$ with $2n$ legs paired into $n$ modes, the
ladder operators satisfy the CAR.
\end{theorem}

\begin{proof}[Proof sketch]
Consider two Majorana operators $S(\ell_a)$ and $S(\ell_b)$.

\emph{Case 1: paths diverge at some ancestor.}
The two root-to-leg paths share a common prefix and diverge at some
node~$v$, where $S(\ell_a)$ follows one child (label $\sigma_a$) and
$S(\ell_b)$ follows another (label $\sigma_b$, with $\sigma_a \neq
\sigma_b$).  At node~$v$: $\sigma_a \sigma_b = -\sigma_b \sigma_a$
(anticommutation of distinct Paulis).  At all shared prefix
nodes: same labels, so they commute.  Below the branch: one operator
has identity where the other is non-trivial.  Total anticommuting
positions: exactly~1 (odd), so $\{S(\ell_a), S(\ell_b)\} = 0$.

\emph{Case 2: one path is a prefix of the other.}
This cannot occur because legs are terminals (they don't continue
to child nodes), so two distinct legs always diverge.

\emph{Case 3: same leg ($\ell_a = \ell_b$).}
$S(\ell)^2 = I^{\otimes n}$ (each Pauli squares to identity),
so $\{S(\ell), S(\ell)\} = 2I$.

In all cases, the Majorana anticommutation
relations~\eqref{eq:majorana-car} are satisfied.  The formal
inductive proof is given by Jiang et al.~\cite{jiang2020}.
\end{proof}

The elegance of Construction~B is that correctness follows from the
\emph{tree structure alone}: distinct paths through a branching tree
necessarily anticommute at exactly one node.  No constraints on tree
shape, depth, or labelling are needed.

\subsection{Recovery of known encodings}
\label{sec:recovery}

All five standard encodings arise from specific tree topologies:

\begin{table}[t]
\centering
\caption{Standard encodings as tree topologies.}
\label{tab:tree-encodings}
\begin{tabular}{@{}lll@{}}
\toprule
Tree topology & Encoding & Max Pauli weight \\
\midrule
Star (depth 1, chain ordering) & Jordan--Wigner & $O(n)$ \\
Star (depth 1, reversed) & Parity & $O(n)$ \\
Fenwick tree & Bravyi--Kitaev & $O(\log_2 n)$ \\
Balanced binary tree & Binary tree & $O(\log_2 n)$ \\
Balanced ternary tree & Ternary tree & $O(\log_3 n)$ \\
\bottomrule
\end{tabular}
\end{table}

The star-tree topologies (JW, Parity) have maximum depth~1, giving
weight up to $2(1) + 1 = 3$ \emph{per Majorana operator} --- but the
star has $n$ legs radiating from a single root, and the root qubit
participates in \emph{every} Majorana operator, leading to $O(n)$
weight in aggregate.  Deeper trees distribute the work across more
internal nodes, achieving logarithmic weight.

\subsection{Pauli-weight bounds}
\label{sec:weight-bounds}

\begin{theorem}[Weight bound]
\label{thm:weight-bound}
For an encoding tree $T$ with depth $d = \depth{T}$, every Majorana
operator has Pauli weight at most $d + 1$.  Every encoded ladder
operator has Pauli weight at most $2d + 1$.
\end{theorem}

\begin{proof}
Each Majorana string has weight equal to its root-to-leg path length,
at most $d + 1$.  A ladder operator is
$\tfrac{1}{2}(S(\ell^{(c)}) \pm i\,S(\ell^{(d)}))$; in the worst
case, the two Majorana strings share only the root, giving weight at
most $2d + 1$.
\end{proof}

\begin{corollary}[Optimal asymptotic weight]
\label{cor:optimal-weight}
Since each node has at most 3 children (one per Pauli label), the
minimum depth for $2n$ legs is $d \ge \lceil\log_3(2n)\rceil$.
The balanced ternary tree achieves this bound, giving optimal weight
$O(\log_3 n)$ among fixed-arity trees.
\end{corollary}

This optimality result, established by Jiang et al.~\cite{jiang2020},
is notable: the provably optimal encoding uses a balanced ternary tree
--- a tree that is \emph{not} a star, and therefore lies outside the
domain of the index-set construction (\cref{sec:star-tree}).
